% 第5章 深度学习与序列建模
\section{深度学习与序列建模}\label{sec:deep}

如果说监督学习那章讨论的是“在横截面上找规律”，这一章讨论的是“在时间轴上找规律”。两者看似只差一个维度，实际却把问题换了个性质：横截面预测关心的是同一时点上资产之间的相对高低，时序预测关心的是同一资产在不同时刻的演化。后者的难点，一半来自序列本身，一半来自金融时间序列那点特殊的坏脾气。

\subsection{问题设定：预测精度不等于交易绩效}

序列建模的目标通常是预测未来的价格、收益或波动率。这里必须先把一件事分清楚，否则后面全乱：\textbf{预测精度（如 RMSE、方向准确率）与交易绩效不是一回事。}一个能把价格拟合得很准的模型，可能因为预测的方向变化太频繁、幅度太小、或者领先/滞后关系不对，而在扣掉成本后赚不到钱。大量“LSTM 预测股价”的论文只报告了精度指标，却从未把它接到一个带成本的交易回测上——这类工作在本文里只作背景，不计入可交易证据。

\subsection{架构谱系}

\textbf{循环网络}。Fischer 与 Krauss（2018）\cite{fischerkrauss2018} 用 LSTM 预测标普 500 成分股的次日方向，报告了约 0.46\% 量级的日度方向准确率提升与对应的回测结果，是 LSTM 进入金融预测的代表作。它的意义不只在于结果，更在于它示范了一种做法：\textbf{把预测接到一个明确的交易规则上，并报告成本敏感性}，而不是停在精度上。

\textbf{卷积与混合结构}。一维卷积（TCN）与卷积--循环混合结构在捕捉局部模式上有优势，常被用作对循环结构的替代或补充。

\textbf{注意力与 Transformer}。Vaswani 等（2017）\cite{vaswani2017} 的注意力机制进入时序领域后，出现了两类工作：一类是通用的长序列预测架构（Informer\cite{zhou2021informer}、Autoformer\cite{wu2021autoformer} 等），另一类是为多变量、带已知协变量的预测设计的结构，如 Temporal Fusion Transformer（Lim 等 2021\cite{lim2021tft}），后者因其可解释的变量选择与门控机制，在含外生变量的金融场景里被较多采用。

这里要专门提一个反证。Zeng 等（2023）\cite{zeng2023dlinear} 指出，在多个通用时序预测基准上，一个极其简单的线性模型（DLinear）就能追平甚至超过上述复杂 Transformer，说明“注意力带来的收益”有时被高估了。这个结论针对的是通用时序基准，不能直接外推到金融，但它给的警示是普适的：\textbf{在比较模型时，一个强线性基线不可缺少；既然低信噪比会放大复杂度带来的过拟合，金融场景里尤其如此。}

\textbf{时序基础模型}。2023 年后出现了一批在跨领域大规模时序上预训练的“基础模型”，如 Chronos（Ansari 等 2024\cite{ansari2024chronos}）、TimesFM（Das 等 2024\cite{das2024timesfm}）、Lag-Llama（Rasul 等 2023\cite{rasul2023lagllama}）。它们的目标是零样本/少样本预测。金融是这类模型最诱人、也最危险的应用场：诱人在于数据量大、任务多；危险在于金融序列的非平稳与对抗性，与这些模型预训练时的“平稳、可重复”假设相冲突。目前公开的、在严格样本外协议下把时序基础模型接到可交易策略上的证据仍然稀少，本文将其列为\textbf{证据强度低--中、需要独立复现}的方向。

\subsection{一个常被忽略的落地点：订单簿与微观结构}

序列建模在金融里最干净的落地场景，其实不是日频价格预测，而是\textbf{订单簿建模}。订单簿数据的样本量以十亿计，信噪比虽不高但相对稳定，且预测目标（下一笔成交方向、未来的短时价格漂移）与执行决策直接挂钩。Zhang、Zohren 与 Roberts（2019）\cite{zhang2019deeplob} 的 DeepLOB 用卷积加 LSTM 处理限价订单簿，Sirignano（2019）\cite{sirignano2019lob} 提出面向订单簿的空间网络结构，Tsantekidis 等（2017）\cite{tsantekidis2017} 用卷积网络从订单簿预测价格方向，这些工作之所以相对可信，正是因为它们的数据与评测协议比日频选股要干净得多。本文把订单簿建模放在这一族，是因为它的方法论贡献（如何把高维、非平稳的微观结构序列喂给网络）对其他环节也有借鉴价值；而把它与“用 LSTM 预测股价”区分开，是因为二者的可复现性不在一个量级上。

\subsection{落地要素}

\begin{itemize}
  \item \textbf{标签构造决定一切}。用未来 1 日、5 日还是 20 日收益作标签，等价于选择了不同的预测尺度，动量和反转会跟着翻转。报告结果时必须写清楚。
  \item \textbf{标准化要按时间在线进行}。用全样本均值方差做标准化，是把未来信息泄漏进训练集的典型错误。
  \item \textbf{验证集要用滚动/前向切分}。随机切分在时序问题上等同于作弊。
  \item \textbf{对齐交易时点}。信号在 $t$ 收盘算出、却假设在 $t$ 开盘成交，是最常见的隐性前瞻偏差。
\end{itemize}

\subsection{复现状态与失效边界}

这一族的开源实现很多（各类 LSTM/Transformer 选股仓库），但真正带完整成本模型、样本外协议与独立复核的很少。DeepLOB 因为数据来源（FI-2010\cite{ntakaris2018fi2010} 等公开订单簿数据集）与代码相对透明，复现状态较好；而多数“深度模型选股”论文的复现状态为\textbf{部分复现或不可复现}。

失效边界可归纳为三条：\textbf{非平稳}（模型学到的时序模式在制度切换后失效）、\textbf{过参数化}（在低信噪比下把噪声当信号）、以及\textbf{目标错位}（优化预测精度而非交易绩效）。本章要记住的一句话是：\textbf{在金融序列上，模型能不能“预测得准”和能不能“赚到钱”是两件事，而绝大多数标题党式的深度学习论文，恰恰混淆了这两件事。}
